% /root/monodromy, 19 November 2026; independently audited by /root.
\subsection{Many fixed covectors for ordinary abelian ambient varieties}
\label{subsec:ordinary-abelian-covectors}
\status{Proved; independent parent proof and source audit passed.}

\begin{proposition}\label{prop:ordinary-abelian-fixed-quotient}
Let $A$ be an ordinary abelian variety of dimension $n\geq2$
over a perfect field of characteristic two, with fixed very
ample $L=\mathcal O_A(1)$.  For every sufficiently large even
degree $d=2t$, the geometric monodromy of $\mathcal V_d$
acts trivially on a quotient of dimension $2^n-1$.
More precisely, write $F:A\to A^{(2)}$ and, after a field
extension splitting $\ker F$, write
\[
 F_*\mathcal O_A=\bigoplus_{\chi\in H}\mathcal N_\chi,
 \qquad H=\operatorname{Hom}(\ker F,\mathbf G_m)
                 \simeq(\mathbf Z/2)^n.
\]
Set
\[
 R=\max_{\chi\in H}
         \operatorname{reg}_{L^{(2)}}(\mathcal N_\chi).
\]
It suffices that $t\geq\max\{1,R\}$ and
$2t\geq\max\{d_{\rm coh},\operatorname{reg}(\mathcal I_A)-1\}$.
After the same splitting extension, the arithmetic image
also fixes this quotient pointwise.
\end{proposition}

\begin{proof}
The character decomposition is the decomposition of the
algebra of a torsor under the diagonalizable group
$\ker F\simeq\mu_2^n$.  In particular,
$\mathcal N_0=\mathcal O_{A^{(2)}}$ and multiplication gives
$\mathcal N_\chi\otimes\mathcal N_\psi\simeq
\mathcal N_{\chi+\psi}$.
The line bundles are pairwise nonisomorphic two-torsion
bundles in $\Pic^0(A^{(2)})$.
This decomposition, including all the two-torsion line
bundles, is also stated and proved in
\cite[Lemma~4.5, p.~11]{SannaiTanaka2016}.

Choose the unique nonzero invariant top differential
$\omega$ normalized by $C(\omega)=\omega$, where $C$
is absolute Cartier.  Such a normalization exists because
$A$ is ordinary and the characteristic is two: if
$C(\omega_0)=c\omega_0$ with $c\ne0$, take
$\omega=c^2\omega_0$.
The coefficient trace $F_*\mathcal O_A\to\mathcal O_{A^{(2)}}$
defined by this differential is the projection to the
trivial character.  Indeed its restriction to every
nontrivial character line is zero, since a nontrivial
degree-zero line bundle has no global sections; on
$\mathcal N_0$ the normalization makes it the identity.

Because $L^{2t}=F^*(L^{(2)t})$, the projection formula
gives a direct sum decomposition
\begin{equation}\label{eq:abelian-section-character-decomposition}
 W:=H^0(A,L^{2t})
    =\bigoplus_{\chi\in H}
       H^0(A^{(2)},L^{(2)t}\otimes\mathcal N_\chi).
\end{equation}
Write $s=\sum_\chi s_\chi$ for the corresponding
character components, regarded again as sections on $A$.
Each summand in \eqref{eq:abelian-section-character-decomposition}
is nonzero under the stated regularity bound.  Thus for
generic $s$ all its $2^n$ components are nonzero, and
these components are linearly independent.

We claim that each
\[
 w_\chi(s)=s_\chi\omega\in H^0(A,\omega_A L^{2t})
\]
is a Frobenius-equivariant covector for $F_s=sF_A$.
The twisted Cartier operator dual to $F_s$ sends it to
$C(s s_\chi\omega)$.  In the sum
$s s_\chi=\sum_\psi s_\psi s_\chi$, every term with
$\psi\ne\chi$ has nontrivial character and is killed
by Cartier, whereas
\[
 C(s_\chi^2\omega)=s_\chi C(\omega)=s_\chi\omega.
\]
Consequently $C_s(w_\chi)=w_\chi$.
The character projections are linear in $s$, so these
are global covectors on the tautologically twisted
coherent bundle, just as in
Theorem~\ref{thm:general-adjunction-covector}.

By Serre duality the dual of the relative coherent
fibre is $W/k s$, since $H^0(A,\omega_A)=k\omega$.
The $2^n$ independent vectors $s_\chi$ span a subspace
whose intersection with $ks$ is exactly $ks$.
Their classes in $W/ks$ therefore have dimension
$2^n-1$, with the single relation
\begin{equation}\label{eq:abelian-covector-single-relation}
 \sum_{\chi\in H}w_\chi=0
       \quad\hbox{on the relative coherent fibre}.
\end{equation}
They kill the nilpotent part by Frobenius equivariance,
so remain independent on the stable part.  Lang descent
turns them into $2^n-1$ independent $\mathbf F_2$-valued
covectors; independence can be checked after tensoring
with the algebraic closure.  Their rank is constant on
the connected maximal-rank open.  The ambient restriction
bound in the statement permits pullback from the complete
linear system to the prescribed hypersurface parameter
space.  This proves the asserted fixed quotient.
\end{proof}

\begin{remark}\label{rem:abelian-obstruction-not-exact-group}
Already for an ordinary abelian surface there is a fixed
quotient of dimension three.  Thus even the proposed
replacement ``full stabilizer of one adjunction covector''
would be too large for general ambient varieties.
The proposition gives a lower bound for the full fixed
covector space; it does not assert that no further
invariants occur, nor determine the whole monodromy group.
\end{remark}

\begin{example}\label{ex:ordinary-abelian-surface-explicit}
Here is a completely specified ambient surface, with an executed
ordinary-hypersurface calculation.  Over $\mathbf F_2$, let
\[
 E:\quad Y^2Z+XYZ+X^3+Z^3=0,\qquad
 A=E\times E,\qquad L=\mathcal O_E(3O)\boxtimes\mathcal O_E(3O),
\]
where $O=[0:1:0]$, and embed $A$ in $\mathbf P^8$ by the
Segre embedding of these two plane cubics.  The derivatives of
the displayed cubic have no common projective zero on $E$;
its four $\mathbf F_2$-points are
$[0:1:1],[1:0:1],[1:1:1],O$.
Thus $E$ is smooth and ordinary (its Frobenius trace is $-1$).
The degree-one section
\begin{equation}\label{eq:abelian-degree-one-example}
 Z_1Z_2+Z_1Y_2+X_1Y_2+Y_1X_2=0
\end{equation}
cuts out a smooth ordinary genus-ten curve on $A$.

The exact programs \path{ordinary_abelian_surface.py} and
\path{ordinary_abelian_degree_one.py}, with their retained
JSON outputs, verify this as follows.  On $E\setminus\{O\}$
write $x=X/Z,y=Y/Z$ and $\omega=dx/x$.
The coefficient Cartier map satisfies
\[
 T(x^{2a})=x^a,\quad T(x^{2a+1})=0,\quad
 T(x^{2a}y)=x^{a+1},\quad T(x^{2a+1}y)=x^a(y+1).
\]
These formulas follow directly from
$((y+1)/x)^2=y/x+x$ and the decomposition
$k(E)=k(E)^2\oplus xk(E)^2$.
Multiplication uses $y^2=x^3+xy+1$.
On the tensor basis $(1,x,y)\otimes(1,x,y)$ the Cartier
map $g\mapsto T_1T_2(sg)$ has the following nine columns,
encoded as binary integers from the first basis vector:
\[
 (3,21,46,16,170,357,154,493,116).
\]
Its rank over $\mathbf F_2$ is nine.  Quotienting the fixed
line $ks$ gives relative rank eight.  Together with the
ordinary ambient $H^1(A,\mathcal O_A)$ of rank two, this
proves $p$-rank ten.

Smoothness is checked on all four product charts
$Z_i\ne0$ or $Y_i\ne0$.
The invariant tangent vector field has components
$(x,x^2+y)$ on the first chart and
$(1+u+v^2,u^2+v)$ on $u=X/Y,v=Z/Y$.
On each product chart the ideal generated by the two
elliptic equations, the section, and its two tangent
derivatives has Groebner basis $[1]$ over $\mathbf F_2$.
These charts cover $A$ and the vector fields never vanish,
so these certificates prove smoothness over the algebraic
closure.  The direct count is eight rational points,
consistent with the computed Cartier trace zero.

It follows that the generic degree-$d$ section on this
\emph{fixed} embedded surface is ordinary for every $d\geq1$.
Indeed choose $d$ general ordinary members of $|L|$;
they meet transversely, with no triple intersections.
Each has genus ten and each pair meets in $L^2=18$ points.
Their connected nodal union has
\[
 h^1(\mathcal O)=10d+18\binom d2-d+1=9d^2+1.
\]
The normalization sequence expresses this cohomology
as an extension of the ordinary component cohomology
by the Frobenius-bijective graph cohomology.  Its
Frobenius is therefore bijective.  Openness of that
condition and the existence of smooth members give
ordinary generic smooth sections.  Thus
\begin{equation}\label{eq:abelian-example-generic-rank}
 r_d=9d^2-1\qquad\text{for every }d\geq1.
\end{equation}
The ambient restriction map is surjective in every
degree: restriction of plane forms to each cubic is
surjective, and Segre monomials span all forms of
bidegree $(d,d)$.

For every even $d\geq2$, the four character summands
have nonzero spaces of sections (their factors on each
elliptic curve have degree $3d/2$), so the fixed quotient
of Proposition~\ref{prop:ordinary-abelian-fixed-quotient}
has dimension three with no additional degree bound.
The degree-two program also gives an explicit smooth
ordinary genus-37 section: its 36-dimensional Cartier
matrix has rank 36; the relative rank is 35, and its
four character covectors have rank four before taking
the single relation and rank three afterwards.
All defining coefficients, matrix columns, covectors,
chart certificates and the count of seven rational
points are retained in its JSON output.
\end{example}
